\scanpage{8}
% archive foliation: 7
\unsure{donnée} par la condition, (pour $\forall z\in C(j)$)
\[
\Sigma(z)=\begin{cases}
\text{a)}\ R(z)\text{ si }\deleted{\unsure{\nexists}\text{ partie }(x_0,x_1,\ldots,x_n)\text{ avec }x_0\in\Gamma,\ x_n=z}\ \deleted{\unsure{\nexists}x\in\Gamma,\ \unsure{\text{avec}}\ z\unsure{\succ|\geq}x,\ \text{ou si }z\notin G(j)}\ \inserted{z\notin\Gamma}\\
\text{b)}\ R(z)\cap G(j)\text{ si }\deleted{\exists x\in\Gamma\text{ avec }z\unsure{\succ|\geq}x,\text{ et }z\in G(j)}\ \inserted{z\in\Gamma}
\end{cases}
\]
\deleted{\unsure{bien} (donc $z\in G(j)$)}

\uline{Dém} Prouvons que $\Sigma(z)$ est une stratégie gagnante pour $\forall x\in\Gamma$. C'est une stratégie i.e. $R(z)\ne\varnothing\Longrightarrow\Sigma(z)\ne\varnothing$. Il suffit de regarder le cas \deleted{b) $\exists x\in\Gamma$, avec $z\unsure{\succ|\geq}x$, et $z\in G(j)$}\inserted{$z\in\Gamma$} (sinon $\Sigma(z)=R(z)\ne\varnothing$), alors $\Sigma(z)=R(z)\cap G(j)$, et par le $1^\circ)$ de la prop. on a (puisque $z\notin C_0$ donc $z\notin G_0(j)$) $R(z)\cap G(j)\ne\varnothing$ OK

Prouvons que $\Sigma(z)$ est gagnante pour $\forall x\in\Gamma$. Soit donc une partie \inserted{$(x_0=x,x_1,\ldots,x_{\unsure{N|n}})$} compatible avec $\Sigma$ et maximale, \deleted{terminant en}\inserted{partant de} $x$, et \unsure{d'aboutissement} : $x_{\unsure{N|n}}$, prouvons $x_{\unsure{N|n}}\in G(j)$ (Alors on aura, $x_{\unsure{N|n}}\in G_0(j)$). Par hyp $x_0\in G(j)$, on va prouver par rec.
\[
x_i\in G(j)\Longrightarrow x_{i+1}\in G(j).\qquad(0\leq i\leq N-1)
\]
\uline{Si} $x_i\in\overline C(j)$, cela résulte de Prop., $2^\circ$. Si $x_i\in C(j)$, cela résulte de la définition de $\Sigma$, qui implique que $\Sigma(x_i)=R(z)\cap G(j)$, donc \inserted{comme} $x_{i+1}\in\Sigma(x_i)$ \deleted{\unsure{et z}} $x_{i+1}\in G(j)$. O.K.

Prouvons enfin que parmi les stratégies $\Sigma$
